\SourceLeaf{294}%
\[
\begin{aligned}
t&=t^0(1-\tfrac{1}{6}\rho\rho rr-\tfrac{1}{3}\rho\rho t^0t^0)\\
\lambda&=\lambda^0(1-\tfrac{1}{6}\rho\rho s^0s^0-\tfrac{1}{3}\rho\rho t^0t^0)\\
\sigma&=\tfrac{1}{2}\rho vt^0(1-\tfrac{1}{12}\rho\rho rr-\tfrac{1}{4}\rho\rho s^0s^0-\tfrac{1}{4}\rho\rho t^0t^0)\\
\tau&=\tfrac{1}{2}\rho vs^0(1+\tfrac{5}{12}\rho\rho rr-\tfrac{1}{2}\rho\rho s^0s^0)
\end{aligned}
\]
für \(t\) und \(\lambda\) bis zur fünften, für \(\sigma\) und \(\tau\) bis zur sechsten Ordnung, jeweils ausschließlich. Noch bequemer und ebenso genau ist es, gleich die Logarithmen zu verwenden. Bezeichnet man zur Abkürzung das Produkt der Größe \(\tfrac{1}{12}\rho\rho\) und des Moduls der Logarithmen nach Briggs mit \(\mu\), so erhalten die Formeln folgende Gestalt:
\[
\begin{aligned}
\log s&=\log s^0+4\mu rr-4\mu s^0s^0\\
\log t&=\log t^0-2\mu rr-4\mu t^0t^0\\
\log\lambda&=\log\lambda^0-2\mu s^0s^0-4\mu t^0t^0\\
\log\sigma&=\log\tfrac{1}{2}\rho vt^0-\mu rr-3\mu s^0s^0-3\mu t^0t^0\\
\log\tau&=\log\tfrac{1}{2}\rho vs^0+5\mu rr-6\mu s^0s^0
\end{aligned}
\]
\indent Diese fünf Formeln bilden zusammen mit den vorangegangenen für \(s^0,t^0,\lambda^0\) eine fünfte Lösungsart. Ihr Kennzeichen ist, dass Näherungswerte für \(s,t,\lambda,\sigma,\tau\) durch kleine, sehr leicht zu berechnende Korrekturen an den Logarithmen zu genauen Werten werden. Die dabei verwendeten konstanten Logarithmen sind
\[
\begin{aligned}
\log\rho&=4{,}6855748668&&(-10)\\
\log\tfrac{1}{2}\rho&=4{,}3845448712&&(-10)\\
\log\mu&=7{,}9297527989&&(-20)
\end{aligned}
\]
oder, wenn die Korrekturen unmittelbar in Einheiten der siebten Dezimalstelle erscheinen sollen,
\[
\log\mu=4{,}9297527989\qquad(-10)
\]
Von diesen Logarithmen werden hier jedoch nur die ersten Ziffern verwendet.
\Article{17}
Die Beziehungen zwischen den Größen \(r,S,S',T,T',\lambda\) lassen sich aber viel einfacher ausdrücken, wenn man vom Mittel der beiden Breiten
